\section{Failure Modes}
\label{sec:failure-modes}

\subsection{One error class for every authentication-class failure}

\pqcrypt{}'s \texttt{decrypt} API returns exactly one error category for every failure where the attacker could be probing for an oracle: \texttt{DecryptionError} (Python, TypeScript, Java) / \texttt{TESSERA\_SEAL\_E\_DECRYPT\_FAILED} (C ABI) / \texttt{Error::Decrypt} (Rust). The categories that map to this single error:

\begin{itemize}
\item Bad AEAD tag.
\item Bad UtC commitment tag.
\item Bad KEM ciphertext (ML-KEM implicit-reject path).
\item AD mismatch.
\item Unknown variant inside a decrypted internal blob (reserved for future suite-version embeds).
\end{itemize}

The categories that map to distinct \emph{structural} errors and may be distinguished safely:

\begin{itemize}
\item Truncated envelope or impossible length field (\texttt{INVALID\_LEN}).
\item Unknown variant byte at offset 6 (\texttt{UNKNOWN\_VARIANT}).
\item Unknown suite byte at offset 7 (\texttt{UNKNOWN\_SUITE}). For non-hybrid-capable language ports (v1.x), suite \code{0x06} is rejected here: at \code{peek\_header} \emph{and} at \code{decrypt}, and ordered \emph{before} the \texttt{BAD\_KEY} suite/key-consistency check below. The ordering is pinned by per-language unit tests in all five non-hybrid ports (TSEAL-00044).
\item Suite/key mismatch (\texttt{BAD\_KEY}). \code{encrypt} rejects when the recipient public key's suite does not match the requested suite; \code{decrypt} rejects when the private key's suite does not match the envelope's header suite. This is structural --- it does not touch secret-correlated bytes --- and distinguishable safely. (Enforced in all eight reference implementations as of 2026-06-10.)
\item Output buffer too small (C ABI only: \texttt{BUF\_TOO\_SMALL}).
\end{itemize}

\subsection{Constant-time discipline}

On every \emph{authentication-class} failure path (bad AEAD tag, bad commit tag, KEM implicit-reject, AD mismatch), the decrypt routine MUST run the synthetic-key schedule --- the HKDF derivations of \cref{sec:kdf-nonce} performed against a synthetic key and then discarded --- as a timing pad, so the failure path performs the same key-derivation work as the success path. The synthetic schedule derives from \texttt{HKDF\_Expand(prk\_zero, "pqcrypt v1 timing pad", 32)} where \texttt{prk\_zero} is a 32-byte all-zero buffer.
% frozen v1 wire tag — not a rebrandable name: "pqcrypt v1 timing pad"
The literal label \texttt{"pqcrypt v1 timing pad"} is a frozen v1 wire tag (\cref{sec:kdf-nonce}) and is preserved verbatim across the rename to \pqcrypt{}.

The Python reference additionally executes a synthetic AEAD verify on the KEM-decaps-failure path. \emph{Structural} errors --- invalid length, unknown variant, unknown suite, \texttt{BAD\_KEY} suite mismatch --- MAY return early without running the synthetic schedule: they are caller errors that are distinguishable by design (see the structural-error list above), and their detection touches no per-key secret-correlated bytes, so an early return discloses nothing beyond the error class itself. The C reference's failure-path comment was aligned to this implemented behavior in TSEAL-00043. (An earlier revision of this section required a synthetic AEAD verify \emph{even on structural failure paths}; no port implemented that stronger wording, and this revision replaces it with the contract the eight implementations actually share.)

\subsection{Timing measurement (automated script; not yet a blocking gate)}

An automated timing-histogram regression ships at \texttt{tools/scripts/timing-histogram.py}
(TSEAL-00048): it measures $N$ good-vs-bad decrypt pairs against the Python reference (default
$N = 20\,000$; ``bad'' flips one bit in the on-wire commitment tag, so both sides execute the full
KEM-decaps + KDF + AEAD path), reports median and p99 deltas, writes a JSON artifact for cross-run
trend tracking, and asserts a soft threshold on the median per-pair delta (warn above 50\,ns,
hard-fail above 500\,ns, with a \texttt{--strict} flag). It runs as a dispatch-only job in the manual
\texttt{tessera-seal.yml} workflow; it is \textbf{not yet a blocking PR gate}. On contended hosts the
single-run noise floor of this statistic is on the order of $\pm 1\,\mu$s (initial 2026-06-10 runs on
a heavily loaded Apple M4 Max measured medians from $-0.65$ to $+1.4\,\mu$s with sign flips across
runs --- noise, not signal); read the artifact as a trend point. No timing dataset and no timing
assertion exist under \texttt{tessera-seal/conformance/vectors/} --- the regression is a script, not
a conformance vector.

The most recent \emph{manual} measurement (Python reference, 2026-05-19, Apple M4 Max, quiet host)
compared 100\,000 decrypt pairs (50\,000 good envelopes, 50\,000 single-bit-flipped) and observed a
median good-vs-bad latency difference below 50\,ns; the procedure is recorded in
\texttt{tessera-seal/docs/security-audit-v1.2.md} (open item \#1 --- the script above addresses the
automation half of that item; promotion to a blocking CI gate remains open). Per-language
measurement environments are recorded under \texttt{tessera-seal/benchmarks/}.

\subsection{Implicit-reject semantics}

The FO transform in ML-KEM-1024 \citep{fips203,hofheinz2017fo} returns a pseudorandom 32-byte value on malformed ciphertext rather than raising an error. \pqcrypt{} preserves this: a bad \code{KEM\_CT} produces an implicit-reject \texttt{ss\_mlkem}, which is then expanded into a useless $K_{\mathrm{aead}}$, which then fails AEAD tag verification, which finally produces the generic \texttt{DecryptionError}. The attacker cannot distinguish ``bad KEM CT'' from ``bad AEAD tag'' from ``bad commit tag'' from ``AD mismatch.''
